\documentclass[margin,line]{res}
\usepackage{comment}
\usepackage{amsmath}
\usepackage{amsfonts}
\usepackage{hyperref}
\hypersetup{colorlinks=true, linkcolor=blue, filecolor=magenta, urlcolor=cyan}
\usepackage{graphicx,calc}
\RequirePackage{fontawesome5}
\usepackage[utf8]{inputenc} % Ajout pour les caractères français
\usepackage{textcomp,gensymb}
%\def\ci#1{\textcircled{\resizebox{.6em}{!}{#1}}}
\topmargin=-0.2in
\oddsidemargin -.5in
\evensidemargin -.5in
\textwidth=6.0in
\textheight=10.0in
\itemsep=0in
\parsep=0in

\newenvironment{list1}{
 \begin{list}{\ding{113}}{%
  \setlength{\itemsep}{0in}
  \setlength{\parsep}{0in} \setlength{\parskip}{0in}
  \setlength{\topsep}{0in} \setlength{\partopsep}{0in} 
  \setlength{\leftmargin}{0.17in}}}{\end{list}}
\newenvironment{list2}{
 \begin{list}{$\bullet$}{%
  \setlength{\itemsep}{0in}
  \setlength{\parsep}{0in} \setlength{\parskip}{0in}
  \setlength{\topsep}{0in} \setlength{\partopsep}{0in} 
  \setlength{\leftmargin}{0.2in}}}{\end{list}}

\pagenumbering{arabic}


\begin{document}


\name{\vspace{-1.9cm}\huge{Roman Korol \href{https://qckorol.ca/fr}{\faExternalLink*}\;\;\href{https://scholar.google.com.ua/citations?hl=fr&user=YAAMduoAAAAJ}{\faGraduationCap}\;\;\href{https://orcid.org/0000-0002-9275-5897}{\faOrcid}\;\;\href{https://github.com/korolgroup}{\faGithub} \vspace*{.1in}}
 \normalfont \normalsize \vspace{0.09cm} \newline \href{mailto:Roman.Korol@USherbrooke.ca}{\faEnvelope}\;\;Roman.Korol@USherbrooke.ca\vspace{0.1cm}}
% \newline \faPhone \;\;+1 626 219 3728\vspace{0.1cm}


\begin{flushright}


\vspace{-5mm}
\hspace{1in} Université de Sherbrooke

\vspace{-5mm}
 Département de chimie

\vspace{-5mm}
\hspace{1.2in}Bureau : D1-3033

\vspace{-5mm}
 \hspace{2in}Tél. : 819-821-8000, poste 61010

\vspace{-5mm}
 \hspace{2in}Sherbrooke, QC J1K 2R1

\vspace{0mm}
\hspace{0.6in}

\vspace{0mm}


\end{flushright}


\begin{resume}

%
\section{\sc Poste \\Actuel}
{\bf Université de Sherbrooke}, Sherbrooke, Québec, Canada\\
\vspace*{-.1in}
\begin{list1}
\item[] Professeur adjoint, 2026--présent
\end{list1}
\section{\sc Formation}
{\bf Université de Rochester}, Rochester, New York\\
\vspace*{-.1in}
\begin{list1}
\item[] Chercheur postdoctoral, 2023--2025
\vspace*{-.1in}
\end{list1}
{\bf California Institute of Technology}, Pasadena, California\\
\vspace*{-.1in}
\begin{list1}
\item[] Doctorat en chimie théorique, 2018--2023
\vspace*{-.1in}
\end{list1}
{\bf Université de Toronto}, Toronto, Canada
\begin{list1}
\item[] B.Sc. (Honours) en chimie avec haute distinction,  2014--2018
\end{list1}

%
%\section{\sc Academic Employment}

%{\bf University of Toronto}, Toronto, Ontario Canada\\
%\vspace*{-.1in}
%\begin{list1}
%\item[] Teaching Assistant, Mathematics, 2017--
%\end{list1}

%{\bf Quest University Canada}, Squamish, British Columbia Canada\\
%\vspace*{-.1in}
%\begin{list1}
%\item[] Professor, Mathematics, 2013--2017
%\end{list1}

%



%%
%\section{\sc Course Coordination}

%\textbf{MAT135/136 Course Coordinator, University of Toronto} \hfill {2017-2018} 

%\begin{list2}
%\item Managed teams of approximately 13 instructors and 45 TAs
%\item Designed learning goals, assessments,  homework, tutorials, and other curricular material 
%\item Administered courses
%\item Conducted classroom observations of instructors
%\end{list2}

\section{\sc Bourses et \\Prix}
\vspace{2mm}
Prix \href{https://www.rochester.edu/college/gradstudies/postdocs/steadman-awards.html}{Steadman} \hfill{2024}

\vspace*{-3mm}

Prix du leadership étudiant \hfill {2024}

\vspace*{-3mm}

Conférence primée Gray-Hill \hfill{2023}

\vspace*{-3mm}

Bourse d'études supérieures Patricia Beckman \hfill {2018}

\vspace*{-3mm}

Bourse commémorative Michael Rebryk \href{http://ucu-building-community.blogspot.ca/2016/03/ucu-adds-new-scholarships-in-2016.html}{}\hfill {2018}

\vspace*{-3mm}

Bourse Ivan Szak en chimie \hfill {2018}

\vspace*{-3mm}

Médaille d'argent du St. Michael's College \hfill {2018} 

\vspace*{-3mm}

Prix d'excellence de l'Université de Toronto \href{https://registrar.utoronto.ca/finances-and-funding/utea/}{}\hfill{2018}

\vspace*{-3mm}

Bourse d'études du St. Michael's College \hfill {2018}

\vspace*{-3mm}

Médaille d'argent de la Société canadienne de chimie \href{https://www.cheminst.ca/awards/student/}{} \hfill{2018} 

\vspace*{-3mm}

Programme de recherche d'été de premier cycle du \href{https://cqiqc.physics.utoronto.ca/cqiqc-programs/undergraduates/}{CQIQC}\hfill{2017}

\vspace*{-3mm}

Prix d'excellence de l'Université de Toronto \href{https://registrar.utoronto.ca/finances-and-funding/utea/}{}\hfill{2017}

\vspace*{-3mm}

Bourse F. E. Beamish en chimie \hfill{2017}

\vspace*{-3mm}

Bourse d'études de la Caisse populaire \href{http://www.buduchnist.com/scholarships}{Buduchnist}\hfill {2017}

\vspace*{-3mm}

Bourse Ivan Szak en chimie \hfill {2017}

\vspace*{-3mm}

Prix Michael Both pour l'engagement exceptionnel envers la danse %\& Contribution to the Desna Ukrainian Dance Company" 
\hfill {2017}

\vspace*{-3mm}

Bourse commémorative John Melady \hfill{2017} 

\vspace*{-3mm}

Bourse d'études C. W. Burton \hfill{2017} 

\vspace*{-3mm}

Bourse de premier cycle commémorative Gollop en chimie \hfill{2017}

\vspace*{-3mm}

Palmarès du doyen \hfill{2017}

\vspace*{-3mm}

Prix d'excellence de l'Université de Toronto \href{https://registrar.utoronto.ca/finances-and-funding/utea/}{}\hfill{2016}


\vspace*{-3mm}

Bourse de recherche \href{https://www.weizmann.ac.il/feinberg/admissions/kupcinet-getz-international-summer-school/about-program-0}{Kupcinet-Getz}\hfill{2015}

\vspace*{-3mm}


Tableau d'honneur de l'Université de Toronto Mississauga \href{https://www.utm.utoronto.ca/cps/honour-roll/2014-15}{} \hfill{2015}

\vspace*{-3mm}

Bourse d'admission Erindale  \hfill{2014} 

\vspace*{-3mm}

Bourse d'études du Président de l'Ukraine \hfill{2014}\\
{\em(décernée annuellement aux quelque 250 lycéens les plus performants sur plus d'un million)}

\vspace*{-3mm}

Étudiant de l'année de Loutsk\hfill{2014}\\
{\em(décerné annuellement au meilleur diplômé du secondaire sur environ 10 000)}
\vspace*{-3mm}

Premier prix à Intel-Eco Ukraine 2014, l'étape nationale de \href{https://student.societyforscience.org/intel-isef}{Intel ISEF} \hfill{2014}

\vspace*{-3mm}

Médaille d'or à l'Olympiade internationale de projets écologiques \hfill{2013}

%%%%%%%%

%%
\end{resume}
\newpage
\pagestyle{headings}
\name{Roman Korol \vspace*{.1in}}
\begin{resume}
%%
\section{\sc Publications}

\begin{enumerate}
\item[14.] Tiwari, V; Korol, R; Franco*, I Robust Purely Optical Signatures of Floquet States in Laser-Dressed Crystals. \textit{Phys. Rev. Lett.} \textbf{2025}, 135 (18), 186901. DOI: \href{https://doi.org/10.1103/5ywx-7dbs}{10.1103/5ywx-7dbs}
\vspace{0.1cm}
\item[13.] Korol*, R; Chen, X; Franco*, I High-Frequency Tails in Spectral Densities. \textit{J. Phys. Chem. A} \textbf{2025}, 129 (15), 3587-3596. DOI: \href{https://doi.org/10.1021/acs.jpca.5c00943}{10.1021/acs.jpca.5c00943}
\vspace{0.1cm}
\item[12.] Turner*, A C; Korol, R; Bill, M; Stolper, D A Stable isotope equilibria in the dihydrogen-water-methane-ethane-propane system. Part 2: Experimental determination of hydrogen isotopic equilibrium for ethane-H$_{2}$ from 30 to 200 °C and propane-H$_{2}$ from 75 to 200 °C. \textit{Geochim. Cosmochim. Acta} \textbf{2025}, 396, 91-106. DOI: \href{https://doi.org/10.1016/j.gca.2025.02.033}{10.1016/j.gca.2025.02.033}
\vspace{0.1cm}
\item[11.] Korol*, R; Turner, A C; Nandi, A; Bowman, J M; Goddard, W A; Stolper, D A Stable isotope equilibria in the dihydrogen-water-methane-ethane-propane system. Part 1: Path-integral calculations with CCSD(T) quality potentials. \textit{Geochim. Cosmochim. Acta} \textbf{2025}, 396, 71-90. DOI: \href{https://doi.org/10.1016/j.gca.2025.02.028}{10.1016/j.gca.2025.02.028}
\vspace{0.1cm}
\item[10.] Turner, A C; Korol, R; Eldridge, D L; Bill, M; Conrad, M E; Miller, T F; Stolper*, D A Experimental and theoretical determinations of hydrogen isotopic equilibrium in the system CH$_{4}$-H$_{2}$-H$_{2}$O from 3 to 200 °C. \textit{Geochim. Cosmochim. Acta} \textbf{2021}, 314, 223-269. DOI: \href{https://doi.org/10.1016/j.gca.2021.04.026}{10.1016/j.gca.2021.04.026}
\vspace{0.1cm}
\item[9.] Korol, R; Rosa-Raíces, J L; Bou-Rabee, N; Miller*, T F Dimension-free path-integral molecular dynamics without preconditioning. \textit{J. Chem. Phys.} \textbf{2020}, 152 (10), 104102. DOI: \href{https://doi.org/10.1063/1.5134810}{10.1063/1.5134810}
\vspace{0.1cm}
\item[8.] Eldridge, D L; Korol, R; Lloyd, M K; Turner, A C; Webb, M A; Miller, T F; Stolper*, D A Comparison of Experimental vs Theoretical Abundances of $^{1}$$^{3}$CH$_{3}$D and $^{1}$$^{2}$CH$_{2}$D$_{2}$ for Isotopically Equilibrated Systems from 1 to 500 °C. \textit{ACS Earth Space Chem.} \textbf{2019}, 3 (12), 2747-2764. DOI: \href{https://doi.org/10.1021/acsearthspacechem.9b00244}{10.1021/acsearthspacechem.9b00244}
\vspace{0.1cm}
\item[7.] Korol, R; Bou-Rabee, N; Miller*, T F Cayley modification for strongly stable path-integral and ring-polymer molecular dynamics. \textit{J. Chem. Phys.} \textbf{2019}, 151 (12), 124103. DOI: \href{https://doi.org/10.1063/1.5120282}{10.1063/1.5120282}
\vspace{0.1cm}
\item[6.] Korol, R; Segal*, D Machine Learning Prediction of DNA Charge Transport. \textit{J. Phys. Chem. B} \textbf{2019}, 123 (13), 2801-2811. DOI: \href{https://doi.org/10.1021/acs.jpcb.8b12557}{10.1021/acs.jpcb.8b12557}
\vspace{0.1cm}
\item[5.] Korol, R; Segal*, D From Exhaustive Simulations to Key Principles in DNA Nanoelectronics. \textit{J. Phys. Chem. C} \textbf{2018}, 122 (8), 4206-4216. DOI: \href{https://doi.org/10.1021/acs.jpcc.7b12744}{10.1021/acs.jpcc.7b12744}
\vspace{0.1cm}
\item[4.] Korol, R; Kilgour, M; Segal*, D ProbeZT: Simulation of transport coefficients of molecular electronic junctions under environmental effects using Büttiker's probes. \textit{Comput. Phys. Commun.} \textbf{2018}, 224, 396-404. DOI: \href{https://doi.org/10.1016/j.cpc.2017.10.005}{10.1016/j.cpc.2017.10.005}
\vspace{0.1cm}
\item[3.] Korol, R; Kilgour, M; Segal*, D Thermopower of molecular junctions: Tunneling to hopping crossover in DNA. \textit{J. Chem. Phys.} \textbf{2016}, 145 (22), 224702. DOI: \href{https://doi.org/10.1063/1.4971167}{10.1063/1.4971167}
\vspace{0.1cm}
\item[2.] Longobardi, L E; Zatsepin, P; Korol, R; Liu, L; Grimme, S; Stephan*, D W Reactions of Boron-Derived Radicals with Nucleophiles. \textit{J. Am. Chem. Soc.} \textbf{2016}, 139 (1), 426-435. DOI: \href{https://doi.org/10.1021/jacs.6b11190}{10.1021/jacs.6b11190}
\vspace{0.1cm}
\item[1.] Korol*, R V; Yanchuk, O M; Marchuk, O V; Orlov, V F; Moroz, I A; Vyshnevskyi, O A Size Stabilizers in Two-electrode Synthesis of ZnO Nanorods. \textit{Phys. Chem. Solid St.} \textbf{2021}, 22 (2), 380-387. DOI: \href{https://doi.org/10.15330/pcss.22.2.380-387}{10.15330/pcss.22.2.380-387}
\vspace{0.1cm}
\end{enumerate}


\end{resume}

\newpage

\pagestyle{headings}


\name{Roman Korol \vspace*{.1in}}
\begin{resume}

\section{\sc Présentations \\et Prix}

Séminaire CERMM, Université Concordia, Montréal, Québec \hfill{2026}

\begin{list1}
\item[]\underline{Présentation invitée :} ``Effets quantiques dans les systèmes complexes"
\end{list1}

\vspace*{-1mm}

Série ``Lunch and Learn", Département de chimie, Université de Sherbrooke, Québec \hfill{2026}

\begin{list1}
\item[]Présentation : ``Effets quantiques dans les systèmes complexes"
\end{list1}

\vspace*{-1mm}

Séminaire CCQS au département de physique, Université de Rochester, Rochester, New York \hfill{2025}

\begin{list1}
\item[]Présentation : ``How NOT to build control-target gates in semiconductor quantum dots and beyond"
\end{list1}

\vspace*{-1mm}

Série de séminaires d'été sur la science de l'information quantique moléculaire \hfill{2025}

\begin{list1}
\item[]Présentation : ``Opening quantum dynamics on gate-defined quantum dots with RLC circuits"
\end{list1}

\vspace*{-1mm}

Conférence de Rochester \href{https://www.sas.rochester.edu/quantum/cqs/cqs-12.html}{sur la cohérence et la science quantique}, Rochester, New York \hfill{2025}

\begin{list1}
\item[]Affiche : ``Hybrid singlet-triplet qubits enable fault-tolerant single-pulse CNOT gate"
\end{list1}

\vspace*{-1mm}

Conférence américaine \href{https://www.actc2024.org/}{sur la chimie théorique}, Chapel Hill, North Carolina \hfill{2024}

\begin{list1}
\item[]Affiche : ``Analog Simulation of Open Quantum Dynamics"
\end{list1}

\vspace*{-1mm}

Conférence \href{https://www.oxy.edu/academics/student-research/urc/2023-summer-research-program}{Gray-Hill} à l'Occidental College, Los Angeles, California.  \hfill{2023}

\begin{list1}
\item[]\underline{Conférence primée :} ``A window to Earth's past with the help of theoretical chemistry"
\end{list1}

\vspace*{-1mm}

Congrès et exposition \href{https://www.cheminst.ca/conference/canadian-chemistry-conference-and-exhibition/}{canadiens de chimie}, Calgary, Canada\hfill {2022}

\begin{list1}
\item[]Communication orale : ``Accurate quantum statistics from improved path-integrals in imaginary time" 
\end{list1}
\vspace*{-1mm}

Mini-colloque sur la \href{https://www.siminegroup.ca/mms/mms2022}{science moléculaire}, Montreal, Canada\hfill{2022}

\begin{list1}
\item[]Affiche : ``Dimension-free ring-polymer molecular dynamics"
\end{list1}

\vspace*{-1mm}

Réunion de printemps de \href{https://www.morressier.com/o/event/623377e0b300ee00119b311f}{l'ACS}, San Diego, California\hfill{2022}

\begin{list1}
\item[]Affiche : ``Accurate quantum statistics from improved path-integrals in imaginary time"
\end{list1}

\vspace*{-1mm}

Séminaire des sciences géologiques et planétaires à Caltech, Pasadena, California\hfill{2022}

\begin{list1}
\item[]\underline{Conférence invitée :} ``$D$ and $^{13}C$ exchange equilibria using Path-Integral Monte-Carlo"
\end{list1}

\vspace*{-1mm}

Réunion de mécanique statistique de  \href{http://berkeleystatmech.org/}{Berkeley}\hfill{2020}

\begin{list1}
\item[]Affiche : ``Cayley modification for strongly stable path-integral molecular dynamics"
\end{list1}

\vspace*{-1mm}

Conférence CECAM sur l'électronique \href{https://www.cecam.org/workshop-details/cecampsi-k-research-conference-biomolecular-electronics-biomolectro-246}{biomoléculaire}, Madrid, Spain \hfill {2018}

\begin{list1}
\item[]Affiche : ``Principles of Charge Transport in DNA: from extensive simulations to neural networkds"
\end{list1}

\vspace*{-1mm}

28\textsuperscript{e} Symposium \href{http://catc.ca/cstcc2018}{canadien} sur la chimie théorique et computationnelle, Windsor, Canada \hfill {2018}

\begin{list1}
\item[]\underline{Prix de la meilleure affiche :} ``Charge transport in DNA: From comprehensive simulations to key principles"
\end{list1}

\vspace*{-1mm}

100\textsuperscript{e} Congrès \href{https://www.cheminst.ca/magazine/article/csc-celebrates-100-years-of-chemistry-in-toronto/}{canadien de chimie}, Toronto, Canada \hfill {2017}

\begin{list1}
\item[]\underline{Prix de la meilleure affiche :} ``Tunneling to Hopping Crossover in Thermopower of DNA Molecular Junctions" 
\end{list1}

Symposium de \href{www.chembiophys.ca}{biophysique chimique}, Toronto, Canada\hfill{2017}

\begin{list1}
\item[]Communication orale : ``DNA Molecular Junctions:  Tunneling to Hopping Crossover"
\end{list1}

\vspace*{-1mm}

\href{http://scp.uwaterloo.ca}{33\textsuperscript{rd} Symposium} de physique chimique, Waterloo, Canada\hfill{2017}

\begin{list1}
\item[]Communication orale : ``Probing mechanisms of charge transport in DNA with Landauer-B\"uttiker formalism"
\end{list1}

\vspace*{-1mm}

45e Conférence de chimie des étudiants de premier cycle du Sud de \href{https://sites.google.com/view/souscc52/home}{l'Ontario}, Toronto, Canada \hfill{2017}

\begin{list1}
\item[]\underline{1\textsuperscript{er} prix (présentation) :}  ``Tunneling to Hopping Crossover in DNA \& DNA-like molecular junctions"
\end{list1}

\section{\sc Service}

Comité des études supérieures du département (CESD) \hfill{2026--présent}

\vspace*{-1mm}

Comités de suivi d'étudiants (3 étudiants) \hfill{2026--présent}

\vspace*{-1mm}

Représentant étudiant au comité consultatif du département de chimie \hfill{2016--2017}

\vspace*{-1mm}

Représentant de 2\textsuperscript{e} année à \href{https://csu.sa.utoronto.ca/events-2013-2014/previous-executive-teams/exceutive-team-2015-2016/}{l'association des étudiants en chimie} \hfill{2016--2017}

\vspace*{-1mm}

Membre du conseil d'administration du groupe étudiant {\em Chemistry Connections}\hfill {2015--2016}

\section{\sc Adhésions}

\href{https://www.cheminst.ca/about/about-csc/}{Société canadienne de chimie} (SCC)\\
\href{https://www.concordia.ca/research/molecular-modeling.html}{Centre de recherche en modélisation moléculaire} (CERMM)\\
\href{https://www.usherbrooke.ca/recherche/fr/udes/regroupements/centres}{Centre d'études des matériaux avancés pour les piles et systèmes de stockage d'énergie renouvelable de l'Université de Sherbrooke} (CÉMAPROVUS)

\end{resume}
\newpage

\pagestyle{headings}


\name{Roman Korol \vspace*{.1in}}
\begin{resume}
\section{\sc Bénévolat et \\Initiatives Communautaires}
Direction et coordination de l'envoi de fournitures humanitaires en Ukraine \href{https://www.caltech.edu/about/news/organizing-aid-to-his-native-ukraine}{\faExternalLink}\hfill{2022--2023}

\begin{list1}
\item[]depuis le campus de Caltech et au-delà
\end{list1}

\vspace*{-1mm}

Bénévole pour l'organisation à but non lucratif \href{https://novaukraine.org/}{Nova Ukraine}, Stanford, California \hfill{2021--2023}

\begin{list1}
\item[]Développement web, établissement et coordination du partenariat avec {\em Teach for Ukraine}
\end{list1}

\vspace{-0.3cm}

Bénévole pour l'organisation à but non lucratif \href{https://teachforukraine.org/en/}{Teach for Ukraine}, Kyiv, Ukraine \hfill{2021--2022}

\begin{list1}
\item[]Recrutement et entrevue des enseignants candidats à l'étape de l'entrevue à distance
\end{list1}

\vspace{-0.3cm}

Responsable de \href{https://international.caltech.edu/Orientation/isp}{l'orientation} des étudiants internationaux\hfill{2019, 2020}

\vspace{-0.3cm}

Mentor ("grand frère") pour les nouveaux étudiants aux cycles supérieurs à Caltech\hfill{2019, 2020}

\vspace{-0.3cm}

Bénévole pour le programme de sensibilisation scientifique via \href{https://www.caltechy.org/student-opportunities}{Caltech Y}\hfill{2018--2020}

\vspace{-0.3cm}

Tutorat au secondaire avec l'OBNL \href{https://www.facebook.com/CauseTutoring/?locale=zh_CN&paipv=0&eav=AfapFTOSVzUFy0ircWljNcf8AswCiGwIsBsPs-au8kmxfK99t_d4mxVNzcEVvxME7Jw&_rdr}{CAUSE Tutoring} \hfill{2018--2019}

\vspace{-0.3cm}

Tuteur pour le groupe de tutorat par les pairs de l'Université de Toronto\hfill{2015--2018}

\section{\sc Enseignement}
{\em Développement de cours :} Laboratoires de chimie computationnelle, Chem3 à Caltech\hfill{2022}\\
Accent sur les relations structure-fonction et les dangers des approximations.

\vspace{-0.3cm}

Enseignant au secondaire, Rotman Arts and Science School, Vaughan, Canada \hfill{2020--2022}

\begin{list1}

\item[] Programme académique, Chimie 11\textsuperscript{e} et 12\textsuperscript{e} année, Sciences 10\textsuperscript{e} année. \\
%Adjusted to hybrid curriculum during Covid-19 lockdowns.\\
Un étudiant s'est classé 3\textsuperscript{e} à Vaughan, dans le top 200 au Canada au concours de chimie \href{https://uwaterloo.ca/chemistry/community-outreach/chemistry-high-school-exams}{Avogadro}
\end{list1}

\vspace{-0.3cm}

Entraîneur pour les Olympiades internationales de chimie
\begin{list1}
\item[]Équipe nationale canadienne (4 étudiants) -- 2 médailles de bronze, 1 d'argent \hfill{Printemps 2018}\\
Équipe nationale ukrainienne (1 étudiant) -- médaille de bronze \hfill{Printemps 2014}
\end{list1}

\vspace{-0.3cm}

Tutorat privé en chimie, physique et mathématiques\hfill{2014--2018}

\begin{list1}
\item[]Élèves du secondaire : acceptés à l'University College (R.-U.), l'Université Columbia (É.-U.) et autres
%\\University students: each of them passed the courses I tutored.
\end{list1}

\vspace{-0.3cm}

École d'été de biologie chimique, Loutsk, Ukraine \hfill{Été 2015}

\begin{list1}
\item[]
Conception de problèmes et d'expériences pour aider les lycéens à maîtriser les concepts clés de la chimie
\end{list1}

\section{\sc Évaluation \\par les pairs}
Revues ACS : ACS Physical Chemistry Au\\
Revues APS : Physical Review Letters, Physical Review A, B, E et Research\\
Revues Elsevier : Organic Geochemistry, Chemical Geology\\
Revues RSC : Physical chemistry chemical physics\\
Revues Wiley : Rapid communications in mass spectrometry

\end{resume}
\newpage

\pagestyle{headings}


\name{Roman Korol \vspace*{.1in}}
\begin{resume}

\section{\sc Écoles d'été\\et Ateliers}
%Visiting scholar at the \href{https://sas.rochester.edu/chm/groups/franco/}{Franco group} University of Rochester, New York \hfill{2023}
%
%\vspace*{-1mm}
%
Atelier sur la dynamique en \href{https://meetings.telluridescience.org/meetings/workshop-details?wid=829}{phase condensée} au \href{https://www.telluridescience.org/schools/tellurideschoolontheoreticalchemistry}{TSRC} (Virtuel) \hfill {2020}

\vspace*{-1mm}

École de chimie \href{https://meetings.telluridescience.org/meetings/workshop-details?wid=775}{théorique} au \href{https://www.telluridescience.org/schools/tellurideschoolontheoreticalchemistry}{TSRC}, Telluride, Colorado \hfill {2019}

\vspace*{-1mm}

Institut Weizmann des Sciences, Rehovot, Israël\hfill{2015}\\
\vspace*{-.1in}
\begin{list1}
\item[] Boursier \href{https://www.weizmann.ac.il/feinberg/admissions/kupcinet-getz-international-summer-school/about-program-0}{Kupcinet-Getz} au \href{https://www.weizmann.ac.il/Organic_Chemistry/Rybtchinski/}{laboratoire Rubtchinski}
\end{list1}

\section{\sc Expérience \\Professionnelle \\Antérieure}
Enseignant au secondaire, Rotman Arts and Science School, Vaughan, Canada \hfill{2020--2022}

\vspace*{-1mm}

Assistant de recherche, Département de linguistique, Université de Toronto \hfill{2016--2018}\\
Projet sur la variation et le changement des langues \href{https://ngn.artsci.utoronto.ca/HLVC/0_0_home.php}{d'origine}

\section{\sc Loisirs}

Escalade, haltérophilie \hfill{depuis 2014}\\
Guitare, basse \hfill{depuis 2012}\\
Danse folklorique ukrainienne \hfill{depuis 2004}


\section{\sc Langues}

Maîtrise de l'ukrainien, de l'anglais \textit{\&} du russe
\section{\sc Langages \\Informatiques}

Python, C++, MATLAB, FORTRAN, Mathematica, Bash;\\
Développement web (PHP \& django, HTML, CSS, JS)
\end{resume}
\end{document}